%% LEAN NAMES of the four theorems (file ImprovedExponents/Statements.lean):
%%   \lean{threeSum_paper}    Theorem T2 (the paper's exponents, better parameters)
%%   \lean{threeSum_exact}    Theorem T3 (+ exact count)
%%   \lean{threeSum_full}     Theorem T4 (+ true middle size)
%%   \lean{threeSum_pruned}   Theorem T1 (+ pruned encodings)
%% with exactTriangle_*, minPlus_*, apsp_* for the other problems, the savings
%% allSolved_*_num, and the general theorems allSolved_paper, allSolved_exact,
%% allSolved_full (Pipeline/General.lean), allSolved_pruned (Pipeline/PrunedClaim.lean).

\section{Consequences for 3SUM and APSP}
\label{sec:consequences}

The theorems of this section are statements about programs of the word RAM, in
the sense of Section~\ref{sec:machine}.  Four are for Exact Triangle and its
consequences; they add the ingredients one at a time, and the first adds none.
The fifth is the all-edges route to APSP.

\subsection{From Exact Triangle to the three problems}

\begin{proposition}[\cite{upstream}]\label{prop:transfer}
Let \(0\le\delta\le1\), and suppose that a program solves Exact Triangle on
\(n\) vertices per part with weights up to \(n^{\nu}\) in time
\(O(\nu\,n^{3-\delta}(\log n)^{e})\), uniformly in \(\nu\).  Then
\lean{ThreeSum.SolvedInTime}~\(a\) holds for every rational \(a>2-\delta/2\), and
\lean{MinPlusProduct.SolvedInTime}~\(a\) and \lean{APSP.SolvedInTime}~\(a\) hold for
every rational \(a>3-\delta/3\).
\end{proposition}

These are Theorem~21(a),(b) of the paper as formalized upstream
(\lean{threeSum_of_uniform}, \lean{minPlus_polylog_of_uniform},
\lean{apsp_polylog_of_uniform}).  The hypothesis is obtained from the bound of
Theorem~\ref{thm:total}, which holds from some size on, by solving smaller
instances by brute force (\lean{exactTriangleUniform_of_explicitFrom}).

\subsection{The regime test}
\label{sec:regime}

The upstream solver of the thin product tests whether \(D^{18}\le N\); it runs
the recursion if so and brute force if not.  The test is part of the program
text, and it restricts the ratios at which the program can be used.
\begin{itemize}
\item The recursion must be applicable in the whole regime \(D^{18}\le N\): the
  tile \(\sqrt K N_0\) must fit into \(N\).  For all large \(m\) this holds if
  \(\Apr(c)<18\ln 4\), that is, for \(c<\cFitsMax\ldots\)
  (\lean{exists_fits18}).
\item The instances with \(D\le N^{\varepsilon}\) must lie in the regime:
  \(\varepsilon\le1/18\).  If \(\Afull(c)\ge18\ln 4\), that is, for
  \(c\ge\cFullMin\ldots\), then \(\Rc(\gamma)\le\Rc(0)\le1/18\), and every
  admissible \(\varepsilon\) qualifies (\lean{thinR_le_inv}).
\end{itemize}
Both conditions hold at \(c=\cUncond\) (\lean{basePruned_108_5_lt},
\lean{baseFull_ge_18}), and the numerals of Theorems~\ref{thm:T2}
and~\ref{thm:T3}, which use the program as it is, are obtained there.  The
optimum \(c\approx\cFullApprox\) of the paper's encodings lies just below this
range: there \(\varepsilon\) has to exceed \(1/18\) slightly, while
\(\Apr(c)/\ln4\approx17.6\).

For Theorem~\ref{thm:T4} we therefore add a test with a rational exponent.
Three procedures are appended to the upstream program (\lean{programRS}): a
test \(D^{r}\le N^{s}\) for natural numbers \(r,s\) (\lean{testRSBody}: it
computes \(N^{s}\) by a loop and decides \(D^{r}<N^{s}+1\) with upstream's
routine \lean{powLtBody}, never forming a number above \(N^{s}D+N^{s}+D+r+s\)),
and a copy of the dispatching procedure that calls it.  The analysis of the
solver is carried out once for an abstract regime (\lean{RegimeTest},
\lean{thinClaimR_of_offlineWithin}), of which the upstream test and the new one
are two instances (\lean{regime18}, \lean{regimeRS}).  With the new test the two
conditions above become \(\Apr(c)<(r/s)\ln4\) and \(\varepsilon\le s/r\)
(\lean{exists_fitsRS}, \lean{thinClaimRS_gammaX}).  Such a fraction exists for
every admissible \(\varepsilon\): \(\varepsilon<\Rc(\gamma)\) means
\(1/\varepsilon>(\Afull(c)+\gamma\ln4)/\ln4>\Apr(c)/\ln4\), since
\(\Apr\le\Afull\) (\lean{basePruned_le_baseFull}).  The only remaining
condition on the ratio is \(\varepsilon\le1/16\), needed by the arithmetic of
Theorem~\ref{thm:total}; \(\Afull(c)\ge16\ln4\) for \(c\ge20\) gives it
(\lean{baseFull_ge_16}).

\subsection{The four theorems}

\begin{theorem}[The paper's exponents with better parameters]\label{thm:T2}
With the program, the reduction and the leaf count of the paper, at
\(c=\cUncond\) and suitable rational \(\theta\), \(\varepsilon\) and \(\eta\):
Exact Triangle is solved in time \(O(n^{3-\dTwo})\).  Consequently
\lean{ThreeSum.SolvedInTime} holds with the exponent \(\ThreeSumTwo\), and
\lean{MinPlusProduct.SolvedInTime} and \lean{APSP.SolvedInTime} with the
exponent \(\APSPTwo\).
\end{theorem}

Nothing is new here except the choice of parameters.  Corollary~26 of the paper
rounds \(\varepsilon\) down to \(1/18\) and takes \(c=21\); with
\(\gamma=\min(\gpap(c,\theta),\frac12-q(\theta))\) balanced and
\(\varepsilon\) close to \(\Rc(\gamma)\) the saving
\(\varepsilon\gamma/2\) exceeds \(\dTwo\) (\lean{t2_witness}).  In Lean: the general
statement is \lean{allSolved_paper}, from \lean{thinClaim_gammaOf},
\lean{hostBound_rat} and \lean{explicitFrom_original}; the saving is
\lean{allSolved_paper_num}, and the end statements are
\lean{exactTriangle_paper}, \lean{threeSum_paper}, \lean{minPlus_paper} and
\lean{apsp_paper}.

\begin{theorem}[With the exact count]\label{thm:T3}
With the same program, analyzed by Theorem~\ref{thm:offline}: Exact Triangle is
solved in time \(O(n^{3-\dThree})\).  Consequently
\lean{ThreeSum.SolvedInTime} holds with the exponent \(\ThreeSumThree\), and
\lean{MinPlusProduct.SolvedInTime} and \lean{APSP.SolvedInTime} with the
exponent \(\APSPThree\).
\end{theorem}

In Lean: \lean{allSolved_exact}, from \lean{thinClaim_gammaX} and the same
reduction; the witness is \lean{t3_witness}, the saving
\lean{allSolved_exact_num}, and the end statements are
\lean{exactTriangle_exact}, \lean{threeSum_exact}, \lean{minPlus_exact} and
\lean{apsp_exact}.

\begin{theorem}[With the true middle size]\label{thm:T4}
Let \(c\ge20\) and \(0\le\delta<F_{\Afull(c)}(\Gam(c))\).  With the host of
Theorem~\ref{thm:mid}, the thin product of Theorem~\ref{thm:offline} and the
regime test with a rational exponent, Exact Triangle is solved in time
\(O(n^{3-\delta})\); consequently so is every \(\delta\) below
\(\sup_{c>10}F_{\Afull(c)}(\Gam(c))\in(\SfullLo,\SfullHi]\).  At
\(c=\cUncondOpt\) this holds for \(\delta=\SfullLo\).  Consequently
\lean{ThreeSum.SolvedInTime} holds with the exponent \(\ThreeSumFour\), and
\lean{MinPlusProduct.SolvedInTime} and \lean{APSP.SolvedInTime} with the
exponent \(\APSPFour\).
\end{theorem}

\begin{proof}
By Theorem~\ref{thm:sup} and Proposition~\ref{prop:rational} there is an
admissible rational tuple \((c',\theta',\sigma',\varepsilon')\) near \(c\) with
saving above \(\delta\).  Lower \(\gamma\) below \(\gX(c',\theta')\) and
\(\varepsilon\) to a rational so that the saving stays above \(\delta\)
(\lean{exists_gamma_rat_of_lt}), and choose \(r/s\) between
\(\Apr(c')/\ln4\) and \(1/\varepsilon'\) (Section~\ref{sec:regime}).
Theorem~\ref{thm:offline} with the new test gives the thin product, and
Theorem~\ref{thm:total} with \(\eta=1-1/(2\sigma)\) and
\(d=\varepsilon/(1-\eta)\) slightly lowered gives Exact Triangle.  Its second
condition holds because \(\frac32d\le3\sigma\varepsilon\) and
\(\delta<\varepsilon(1-\sigma)\), so
\(\frac32d+\delta<(1+2\sigma)\varepsilon<3\varepsilon\le\frac3{16}\).  The
supremum over \(c\) is Theorem~\ref{thm:global-full}, and the value at
\(c=\cUncondOpt\) is Proposition~\ref{prop:values}(i).
\end{proof}

In Lean: \lean{allSolved_full} is the statement for every \(c\ge20\),
\lean{allSolved_full_sup} the statement for every saving below the supremum
\lean{SstarFull}, \lean{allSolved_full_num} the saving at \(c=\cUncondOpt\),
and the end statements are \lean{exactTriangle_full}, \lean{threeSum_full},
\lean{minPlus_full} and \lean{apsp_full}.  The five-digit exponents are those
of \(\delta=\dFour\): the optimum improves the saving at \(c=\cUncond\) by less
than \(2\cdot10^{-7}\).

\begin{theorem}[With pruned encodings]\label{thm:T1}
Let \(c\ge21\) and \(0\le\delta<F_{\Apr(c)}(\Gam(c))\).  With the host of
Theorem~\ref{thm:mid} and the thin product of Theorem~\ref{thm:encoder}, Exact
Triangle is solved in time \(O(n^{3-\delta})\); consequently so is every
\(\delta\) below \(\sup_{c>10}F_{\Apr(c)}(\Gam(c))\in(\SstarLo,\SstarHi]\).  At
\(c=\cCond\) this holds for \(\delta=\dOne\).  Consequently
\lean{ThreeSum.SolvedInTime} holds with the exponent \(\ThreeSumOne\), and
\lean{MinPlusProduct.SolvedInTime} and \lean{APSP.SolvedInTime} with the
exponent \(\APSPOne\).
\end{theorem}

The proof is that of Theorem~\ref{thm:T4} with \(\RcP\) in place of \(\Rc\) and
Theorem~\ref{thm:encoder} in place of Theorem~\ref{thm:offline}; the test with
a rational exponent is part of the program \lean{programP}.  For \(c\ge21\) we
have \(\Apr(c)\ge16\ln4\) (\lean{basePruned_ge_16}), so \(\varepsilon\le1/16\)
and \(\frac32d+\delta<3\varepsilon\le0.19\).  The supremum over \(c\) is
Theorem~\ref{thm:global}.  At \(c=\cCond\) the parameters are
\(\Gam>\GamCond\), \(\theta\approx\thetaCond\), \(\sigma\approx\sigmaCond\) and
\(\varepsilon\) just below \(R>\epsCond\) (\lean{Gam_22_bounds},
\lean{thinR_pruned_bounds}, \lean{savingF_pruned_gt}).  In Lean:
\lean{allSolved_pruned} is the statement for every \(c\ge21\) (file
\lean{Pipeline/PrunedClaim.lean}; the form that still takes the thin-product
claim as an argument is \lean{allSolved_pruned_of_claim}, from
\lean{explicitFrom_pruned_of_claim}), \lean{allSolved_pruned_sup} the
statement for every saving below the supremum \lean{SstarPruned},
\lean{allSolved_pruned_num} the saving at \(c=\cCond\), and the end statements
are \lean{exactTriangle_pruned}, \lean{threeSum_pruned}, \lean{minPlus_pruned}
and \lean{apsp_pruned}.  No hypothesis is involved.

\subsection{APSP through all-edges Exact Triangle}
\label{sec:alledges}

Footnote~10 of the paper observes that APSP can keep all of the saving of Exact
Triangle instead of a third: "the proof of Theorem~17 also solves the
all-edges version of Exact Triangle (scanning each pair only until its zero
triangle is found), to which the \(\minplus\)-product reduces on the same \(n\)
vertices".  We carry this out as programs.  The \emph{all-edges version} of
Exact Triangle asks, for every pair \((a,b)\) of the two outer parts, whether
some \(c\) makes \(w(a,b)+w(b,c)+w(a,c)=0\); as a task of the machine
(\lean{aeTask}) it writes the \(n^{2}\) answers to an output array
(\lean{aeFlags}).

\begin{lemma}[Comparisons as equalities]\label{lem:F}
Let \(0\le x,y,u<2^{b}\) be integers and
\(f(\ell)=\lfloor u/2^{\ell}\rfloor-\lfloor x/2^{\ell}\rfloor-\lfloor y/2^{\ell}\rfloor\).
Then
\[
  x+y<u\iff f(0)=1\ \text{ or }\ f(\ell)\in\{2,3\}\text{ for some }\ell\le b .
\]
\end{lemma}

\begin{proof}
\(f(b)=0\), and \(f(\ell)=2f(\ell+1)+e_\ell\), where \(e_\ell\in\{-2,-1,0,1\}\)
is the bit of \(u\) at position \(\ell\) minus those of \(x\) and \(y\).  So
\(f(\ell+1)\ge2\) implies \(f(\ell)\ge2\).  If \(f(\ell)\in\{2,3\}\) for some
\(\ell\), then \(f(0)\ge2\) by descent, and \(f(0)=u-x-y\).  Conversely let
\(f(0)\ge1\).  If \(f(0)\ne1\), let \(\ell\) be the largest index with
\(f(\ell)\ge2\); then \(\ell<b\), \(f(\ell+1)\le1\) and
\(f(\ell)\le2f(\ell+1)+1\le3\).
\end{proof}

In Lean: \lean{add_lt_iff_shiftDiff}.

\begin{proposition}\label{prop:minplus-math}
Let \(A,B\) be \(n\times n\) integer matrices with entries of absolute value at
most \(W\), let \(2W<2^{b}\), and let \(t\) be a matrix of thresholds with
\(-2W\le t_{ij}<2^{b}-2W\).  For \(\ell\le b\) put
\(X_\ell(i,k)=\lfloor(A_{ik}+W)/2^{\ell}\rfloor\),
\(Y_\ell(k,j)=\lfloor(B_{kj}+W)/2^{\ell}\rfloor\) and
\(U_\ell(i,j)=\lfloor(t_{ij}+2W)/2^{\ell}\rfloor\).  Then, for every pair
\((i,j)\), there is a \(k\) with \(A_{ik}+B_{kj}<t_{ij}\) if and only if, for
one of the \(2b+3\) pairs \((\ell,r)\) with \((\ell,r)=(0,1)\) or
\(r\in\{2,3\}\), there is a \(k\) with
\[
  X_\ell(i,k)+Y_\ell(k,j)+\bigl(r-U_\ell(i,j)\bigr)=0 .
\]
Consequently a binary search that runs for all pairs \((i,j)\) in parallel
computes the \(\minplus\)-product of \(A\) and \(B\) from the answers to
\(R(2R+3)\) all-edges Exact Triangle instances on the same \(n\) vertices per
part, where \(R=\lceil\log_2(4W+1)\rceil\).
\end{proposition}

\begin{proof}
Apply Lemma~\ref{lem:F} to \(x=A_{ik}+W\), \(y=B_{kj}+W\) and
\(u=t_{ij}+2W\).  The entry \(C_{ij}\) of the product is the unique integer
\(c\in[-2W,2W]\) such that the answer is negative for the threshold \(c\) and
positive for \(c+1\); a binary search finds it in \(R\) rounds, and a round
asks one matrix of thresholds.
\end{proof}

In Lean: \lean{exists_lt_iff_exists_exactTriangle}, \lean{minPlus_unique},
\lean{search_eq_minPlus} and \lean{search_exactTriangle}.

The reduction of the paper (Theorem~21(b), \cite{VW18}) ends with a binary search of
exactly this kind, over a procedure that answers the threshold questions for
all pairs at once (\lean{pairsTask}, \lean{isHost_mp}), followed by
repeated squaring for APSP (\lean{claim_apspFromMinPlus}).  The factor \(3\) is
lost before that, in a blocking of size \(n^{1/3}\) that reduces the all-pairs
questions to the decision problem.  Two hosts replace the blocking.

\begin{theorem}[All pairs from all-edges Exact Triangle]\label{thm:pairs}
There is a program that, given a procedure solving all-edges Exact Triangle
within \(T(n,U)\) steps, answers the threshold questions of
Proposition~\ref{prop:minplus-math} for all pairs within
\((2b+3)\bigl(T(n,9U)+O(n^{2})\bigr)+O(n^{2})\) steps, where
\(b=\lfloor\log_2 3U\rfloor+1\); in particular within
\(O\bigl((1+\log U)(T(n,9U)+n^{2})\bigr)\).
\end{theorem}

\begin{proof}
With \(W=U\) the hypotheses of Proposition~\ref{prop:minplus-math} hold, and
the weights of the instances are at most \(2^{b}+3\le9U\).  The program writes
\(A+U\), the transpose of \(B+U\) and \(t+2U\), and then runs the \(b\) levels
from the top: at each level it asks the two instances with \(r=2,3\) and takes
the disjunction of their answers into the output, then halves every number
(upstream's pass for \(\lfloor x/2^{\ell}\rfloor\), which avoids division); at
level \(0\) it also asks \(r=1\).  The output is then the matrix of answers.
\end{proof}

In Lean: \lean{isHost_pairsAE} with the time \lean{pairsTimeAE}, from the
bridge \lean{pairFlags_iff} between the two tasks and the invariant
\lean{orFlags_eq_pairFlags}; the bound \lean{pairsTimeAE_le} with the constant
\lean{CA}, and the transfer \lean{pairsFromAllEdges}.

\begin{theorem}[Theorem~\ref{thm:mid} for the all-edges version]\label{thm:mid-ae}
Theorem~\ref{thm:mid} holds for all-edges Exact Triangle in place of Exact
Triangle, with the same bound and another constant \(C\).
\end{theorem}

\begin{proof}
The host of Theorem~17 keeps \(n^{2}\) flags, one per pair \((a,b)\), zero at
the start.  For every query pair that an instance accepts and whose flag is
still \(0\), it scans the piece \(C_k\) of the instance and stores the result
in the flag; pairs whose flag is \(1\) are skipped.  Every pair with a zero
triangle is accepted, and hit, in the instance of the piece of its \(c\), so
the flags are the answers at the end.  A scan either fails, which the paper
charges to a false positive of the prime, or turns a flag from \(0\) to \(1\);
so there are at most \(F(p)+n^{2}\) scans.  The \(n^{2}\) additional scans cost
\(O(n^{2}\sqrt{D'}/g')\), which is within the term \(n^{2}D'g'\); zeroing and
copying the flags costs \(O(n^{2})\).
\end{proof}

In Lean the pure part is \lean{flagsAt_m} (the flags after all instances are
the answers) and \lean{sum_execsAE_le} (the count of the scans); the
procedure that scans the accepted pairs is \lean{scanPairsAEBody}, the brute
force of the small case \lean{bruteAEBody}, the host \lean{ae17_solves}, its
time \lean{hostTimeAE} and \lean{obeysBound17Mid_hostTimeAE}, and the claim
\lean{claim17MidAE_of_pack}.  The statement of Theorem~\ref{thm:mid} is the
same proposition \lean{Claim17Mid} in a time model (\lean{lightModelAE}) whose
reading of "Exact Triangle" is the all-edges task; the deductions of
Sections~\ref{sec:middle} and~\ref{sec:consequences} are made for an arbitrary
model, so they apply unchanged (\lean{MidHostRat}, \lean{midHostRat_ae},
\lean{explicitFrom_full}, \lean{explicitFrom_pruned}).

\begin{theorem}[APSP through all-edges Exact Triangle]\label{thm:apsp-ae}
Let \(c\ge20\) and \(0\le\delta<F_{\Afull(c)}(\Gam(c))\), or let \(\delta\) be
any number below \(\sup_{c>10}F_{\Afull(c)}(\Gam(c))>\SfullLo\).  Then the
\(\minplus\)-product and APSP are solved in time \(O(n^{a})\) for every
rational \(a>3-\delta\); consequently \lean{MinPlusProduct.SolvedInTime} and
\lean{APSP.SolvedInTime} hold with the exponent \(\APSPFn\).  With pruned
encodings the same holds for \(c\ge21\) and \(\delta<F_{\Apr(c)}(\Gam(c))\),
for every \(\delta<\sup_c F_{\Apr(c)}(\Gam(c))>\SstarLo\), and with the
exponent \(\APSPFnOne\).
\end{theorem}

\begin{proof}
Theorems~\ref{thm:T4} and~\ref{thm:T1} with Theorem~\ref{thm:mid-ae} in place
of Theorem~\ref{thm:mid} give all-edges Exact Triangle in time
\(O(n^{3-\delta}\log n)\) for weights up to \(n^{O(1)}\).  Theorem~\ref{thm:pairs}
and the search of Proposition~\ref{prop:minplus-math} (upstream's program,
\(O(\log U)\) rounds at the bound \(6U\)) give the \(\minplus\)-product of
matrices with entries up to \(u\) in time \(O(\log^{2}u\,(T(n,54u)+n^{2}))\),
and repeated squaring gives APSP with \(O(\log n)\) products of matrices with
entries up to \(O(n\,u)\).  Along \(u=n^{\kappa}\) every factor is
\(n^{3-\delta}\) times a polylogarithm.
\end{proof}

In Lean: \lean{minPlusSolved_full}, \lean{minPlusSolved_full_sup},
\lean{minPlusSolved_pruned} and \lean{minPlusSolved_pruned_sup} (the structure
\lean{MinPlusSolved}), through \lean{solvedIn_mp_of_allEdges},
\lean{minPlusInPolylog_of_explicitFrom}, \lean{apspInPolylog_of_explicitFrom}
and \lean{minPlusSolved_of_explicitFrom}; the end statements are
\lean{minPlus_allEdges}, \lean{apsp_allEdges}, and with pruned encodings
\lean{minPlus_allEdges_pruned}, \lean{apsp_allEdges_pruned}, the strongest
statements of the paper for the \(\minplus\)-product and APSP.
